% Generated draft: requires mathematical and source review.
\documentclass[11pt]{article}
\usepackage[margin=25mm]{geometry}
\usepackage{fontspec}
\setmainfont{Latin Modern Roman}
\usepackage{amsmath,amssymb,mathtools,unicode-math}
\setmathfont{Latin Modern Math}
\usepackage{enumitem,needspace,xurl,hyperref}
\setlength{\parindent}{0pt}
\setlength{\parskip}{5pt}
\setlength{\emergencystretch}{3em}
\newcommand{\N}{\mathbb{N}}
\newcommand{\Z}{\mathbb{Z}}
\newcommand{\Q}{\mathbb{Q}}
\newcommand{\R}{\mathbb{R}}
\newcommand{\C}{\mathbb{C}}
\newcommand{\F}{\mathbb{F}}
\newcommand{\A}{\mathbb{A}}
\newcommand{\PP}{\mathbb{P}}
\newcommand{\E}{\mathbb{E}}
\newcommand{\eps}{\varepsilon}
\newcommand{\dd}{\,\mathrm{d}}
\begin{document}
\section*{UnsolvedMath ID: 800878}
\textbf{Source problem code:} catalogue entry
\subsection{Definitions and mathematical statement}
The supplied excerpt leaves the ambient geometric hypotheses and the meaning of its bracket inequality unstated. The following is a conditional formulation in the conventional compact K\"ahler setting; this setting and the interpretation as a total-mass comparison require confirmation against the source context.

Let \(X\) be a compact K\"ahler manifold of complex dimension \(n\geq 1\), and let \(T,S\) be positive closed real \((1,1)\)-currents on \(X\) representing the same cohomology class. Choose a smooth closed real \((1,1)\)-form \(\theta\) in this class and \(\theta\)-plurisubharmonic potentials \(u,v\), so that
\[
T=\theta+\sqrt{-1}\,\partial\bar\partial u,
\qquad
S=\theta+\sqrt{-1}\,\partial\bar\partial v.
\]
Here a \(\theta\)-plurisubharmonic potential is a quasi-plurisubharmonic function whose displayed associated current is positive. Say that \(T\) is more singular than \(S\) if there exists \(C\in\R\) such that \(u\leq v+C\) on \(X\).

Write \(\langle T^n\rangle\) for the non-pluripolar Monge--Amp\`ere measure of \(T\). Locally, if \(T=\sqrt{-1}\,\partial\bar\partial w\) with \(w\) plurisubharmonic, this measure is defined by the increasing limit
\[
\lim_{j\to\infty}
\mathbf{1}_{\{w>-j\}}
\bigl(\sqrt{-1}\,\partial\bar\partial\max\{w,-j\}\bigr)^n,
\]
where the powers of the locally bounded truncations are understood in the Bedford--Taylor sense. Define its total mass by
\[
M(T):=\int_X\langle T^n\rangle.
\]

The total-mass interpretation of the supplied question is: for every such \(X,T,S\), does
\[
[T]=[S]\quad\text{and}\quad u\leq v+C\text{ for some }C\in\R
\quad\Longrightarrow\quad
M(T)\leq M(S)
\]
hold without assuming that either current has small unbounded locus?

The source defines small unbounded locus to mean that local potentials are locally bounded outside a closed complete pluripolar set. It attributes the comparison under this additional assumption for both currents to [BEGZ10, Theorem 1.16], and describes that assumption as apparently technical. The question is attributed there to Boucksom--Eyssidieux--Guedj--Zeriahi.

The excerpt writes \(\langle T^n\rangle\leq\langle S^n\rangle\) without integrals. If those brackets denote measures rather than their total masses or top-degree cohomology classes, the displayed inequality instead means that \(\langle S^n\rangle-\langle T^n\rangle\) is a positive measure. That is a different formulation, and the supplied excerpt alone does not resolve the notation.
\subsection{Short English statement}
Under an explicitly provisional compact K\"ahler and total-mass interpretation, ask whether making a positive closed \((1,1)\)-current more singular, while preserving its cohomology class, can only decrease its non-pluripolar Monge--Amp\`ere mass. The source asks to remove the small-unbounded-locus hypothesis. Its ambient hypotheses and whether its inequality compares masses or measures require confirmation.
\subsection{Sources}
\begin{itemize}
\item Local UnsolvedMath registry, permanent ID \(800878\), exact source problem code \texttt{catalogue entry}; supplied complete source record, upstream revision \texttt{d371ee917ac9647fbb06cd6fc529deec6f3dda82}. The registry labels the item open; this is a supplied status, not an independently verified mathematical or literature assessment.
\item Dinew, Guedj, Zeriahi, \emph{Open Problems in Pluripotential Theory} (2015), Question 9, \url{https://arxiv.org/abs/1511.00705}. The supplied provenance records access on 2026-10-16 and extraction by \texttt{arxiv\_latex}; these are supplied provenance claims, not access performed for this draft.
\item The source attributes the question to Boucksom--Eyssidieux--Guedj--Zeriahi and cites [BEGZ10, Theorem 1.16] for the comparison with small unbounded locus. The supplied record does not include the full bibliographic entry for [BEGZ10].
\end{itemize}
\end{document}
